\chapter{Конструкторская часть}\label{ch:-3}

В данном разделе будут приведены описания алгоритмов умножения матриц, а также модель вычисления трудоемкости алгоритмов.


\section{Описание алгоритмов}\label{sec:-3}
На вход алгоритмов подаются 2 матрицы - $a$ и $b$, на выходе -- матрица $c$, результат умножения матриц $a$ и $b$.
В случае если число столбцов первой матрицы не равно числу строк второй, программа выведет сообщение об ошибке.

На схемах~\ref{scheme:std}---~\ref{scheme:win_opt-2} приведены описания алгоритмов умножения матриц: стандартный алгоритм, алгоритм Винограда, оптимизированный алгоритм Винограда.

\clearpage

\begin{figure}[H]
    \centering
    \includegraphics[scale=0.6]{images/s1}
    \caption{Описание стандартного алгоритма умножения матриц}
    \label{scheme:std}
\end{figure}

\clearpage

\begin{figure}[H]
    \centering
    \includegraphics[scale=0.5]{images/s2}
    \label{scheme:win-1}
\end{figure}

\clearpage

\begin{figure}[H]
    \centering
    \includegraphics[scale=0.8]{images/s3}
    \caption{Описание алгоритма Винограда умножения матриц}
    \label{scheme:win-2}
\end{figure}

\begin{figure}[H]
    \centering
    \includegraphics[scale=0.5]{images/s4}
    \label{scheme:win_opt-1}
\end{figure}

\clearpage

\begin{figure}[H]
    \centering
    \includegraphics[scale=0.8]{images/s5}
    \caption{Описание оптимизированного алгоритма Винограда умножения матриц}
    \label{scheme:win_opt-2}
\end{figure}

Оптимизация алгоритма Винограда по варианту:
\begin{itemize}
    \item двоичный сдвиг вместо умножения на 2;
    \item объединение III и IV частей алгоритма Винограда;
    \item введение декремента при вычислении вспомогательных массивов.
\end{itemize}

\clearpage


\section{Модель вычислений}

Для вычисления трудоемкостей алгоритмов необходимо ввести модель вычислений:

\begin{itemize}
    \item операции $+, -, [], ++, {-}-, ==, !=, <, >, <=, >=$ имеют трудоемкость 1;
    \item операции $*, /, \%, /=,  *=, \%=$ имеют трудоемкость 2;
    \item трудоемкость условного оператора \text{if <condition> then A else B} рассчитывается, как (\ref{eq:ifcond});
    \begin{equation}
        \label{eq:ifcond}
        f_{if} = f_{\text{condition}} +
        \begin{cases}
            f_A, & \text{если условие выполняется,}\\
            f_B, & \text{иначе.}
        \end{cases}
    \end{equation}
    \item трудоемкость цикла \text{for (init; condition; increment) \{ body \}} рассчитывается, как (\ref{eq:forcyc});
    \begin{equation}
        \label{eq:forcyc}
        f_{for} = f_{\text{init}} + f_{\text{condition}} + N\cdot(f_{\text{body}} + f_{\text{condition}} + f_{\text{increment}})
    \end{equation}
    где N -- количество итераций цикла;
    \item трудоемкость вызова функции равна 0.
\end{itemize}


\section{Вычисление трудоемкости алгоритмов}

Далее представлен расчет трудоемкостей алгоритмов умножения матриц на основе модели.
Для расчета представлены такие алгоритмы, как стандартный алгоритм умножения матриц, алгоритм Винограда и оптимизированный алгоритм Винограда.

Проверка входных данных и инициализация результирующей матрицы в расчетах учитываться не будут, так как данные действия не являются самыми трудоемкими и являются общими для всех алгоритмов.

Обозначения размерностей:

\begin{itemize}
    \item $n$ -- число строк первой матрицы;
    \item $m$ -- число столбцов первой матрицы / строк второй матрицы;
    \item $l$ -- число столбцов второй матрицы.н
\end{itemize}

\subsection{Стандартный алгоритм}

Расчет трудоемкости стандартного алгоритма:

\begin{itemize}
    \item трудоемкость внешнего цикла по i, $f_{i} = 2 + n(2 + f_{body_{i}})$;
    \item трудоемкость цикла по j, $f_{j} = 2 + l(2 + f_{body_{j}})$;
    \item трудоемкость цикла по k, $f_{k} = 2 + m(2 + f_{body_{k}})$;
    \item трудоемкость тела цикла по k, $ f_{body_{k}} = 10$,
\end{itemize}

где $f_{body_{i}}$ = $f_{j}$, $f_{body_{j}}$ = $f_{k}$

$f_{std} = 2 + n(4 + l(4 + 12m)) = 12nml + 4mn + 4n + 2 \approx 12nml$

\subsection{Алгоритм Винограда}

Расчет трудоемкости алгоритма Винограда:

\begin{itemize}
    \item создание и инициализация вспомогательных массивов, $n + l$
    \item заполнение первого массива, $f_{1} = 2 + n(2 + 4 + \frac{m}{2}(4 + 14))$
    \item заполнение второго массива, $f_{2} = 2 + l(2 + 4 + \frac{m}{2}(4 + 14))$
    \item заполнение четных размерностей, $f_{3} = 2 + n(4 + l(2 + 7 + 4 + \frac{m}{2}(4 + 26)))$
    \item трудоемкость дополнения если m нечетное $f_{4} = f_{cond}$.
\end{itemize}

\begin{equation}
    \label{eq:cond}
    f_{cond} =
    \begin{cases}
        3, & \text{если $m$ четно (л.с.),}\\
        3 + 2 + n(4 + l(2 + 12)), & \text{иначе (х.с.).}
    \end{cases}
\end{equation}

\begin{equation}
    \label{eq:ans}
    f_{wyn} =
    \begin{cases}
        9 + 11n + 7l + 13ln + 9nm + 9lm + 15nml, & \text{(л.с.),}\\
        11 + 15n + 7l + 27ln + 9nm + 9lm + 15nml, & \text{(х.с.),}
    \end{cases}
\end{equation}

\subsection{Оптимизированный алгоритм Винограда}

Расчет трудоемкости оптимизированного алгоритма Винограда:

\begin{itemize}
    \item создание и инициализация вспомогательных массивов, $n + l$
    \item заполнение первого массива, $f_{1} = 2 + n(2 + 5 + \frac{m}{2}(2 + 12))$
    \item заполнение второго массива, $f_{2} = 2 + l(2 + 5 + \frac{m}{2}(2 + 12))$
    \item заполнение матрицы $c$, $f_{3} = 2 + n(4 + l(2 + 7 + 4 + \frac{m}{2}(4 + 22)) + f_{4})$
    \item трудоемкость дополнения если m нечетное $f_{4} = f_{cond}$.
\end{itemize}

\begin{equation}
    \label{eq:cond-2}
    f_{cond} =
    \begin{cases}
        3, & \text{если $m$ четно (л.с.),}\\
        3 + 12, & \text{иначе (х.с.).}
    \end{cases}
\end{equation}

\begin{equation}
    \label{eq:ans-2}
    f_{optwyn} =
    \begin{cases}
        6 + 12n + 8l + 16ln + 7nm + 7lm + 13nml, & \text{(л.с.),}\\
        6 + 12n + 8l + 28ln + 7nm + 7lm + 13nlm, & \text{(х.с.),}
    \end{cases}
\end{equation}

\section*{Вывод}
В данном разделе были представлены схемы алгоритмов умножения матриц, а также оценены трудоемкости для данных алгоритмов в лучшем и худшем случаях.
Алгоритм Винограда имеет большую трудоемкость по сравнению со стандартным алгоритмом.
Оптимизированный алгоритм имеет меньшую трудоемкость по сравнению с неоптимизированным.

